\documentclass[12pt, a4paper]{article}

\usepackage{xcolor}
\usepackage{geometry}
\usepackage{hyperref}
\usepackage{titlesec}
\usepackage{enumitem}
\usepackage{longtable}
\usepackage{booktabs}
\usepackage{array}
\usepackage{ragged2e}
\usepackage{fancyhdr}

\usepackage[Latin]{xepersian}
\settextfont{Vazirmatn}
\setlatintextfont{Times New Roman}
\setdigitfont{Vazirmatn}

\geometry{margin=2.5cm}

\hypersetup{
    colorlinks=true,
    linkcolor=blue,
    urlcolor=blue
}

\titleformat{\section}{\Large\bfseries\color{blue!60!black}}{\thesection}{1em}{}
\titleformat{\subsection}{\large\bfseries}{\thesubsection}{1em}{}

\setlist{nosep, leftmargin=*}

\pagestyle{fancy}
\fancyhf{}
\fancyhead[R]{\small سند معماری — کلاس یار}
\fancyhead[L]{\small نسخهٔ ۱.۰}
\fancyfoot[C]{\thepage}
\renewcommand{\headrulewidth}{0.4pt}

\linespread{1.4}

\title{
    \vspace{-1cm}
    \textbf{\Huge سند معماری}\\[0.3cm]
    \textbf{\LARGE پروژهٔ کلاس یار}\\[0.5cm]
    \large Architecture Document
}

\author{
    تهیه‌کنندگان: مریم قندی زنجانی، ملیکا شهبازی، آنیسا فروغمند،\\
    مائده رضایی، محمدمهدی ذوالقدر\\
    \vspace{0.5cm}
    دانشگاه آزاد زنجان
}

\date{۲۹/۰۷/۱۴۰۵}

\begin{document}

\maketitle
\thispagestyle{empty}

\vspace{1cm}

\begin{center}
\begin{tabular}{|l|l|}
\hline
\textbf{نام پروژه} & کلاس یار (ClassYar) \\
\hline
\textbf{نسخهٔ سند} & ۱.۰ \\
\hline
\textbf{تاریخ} & ۲۹/۰۷/۱۴۰۵ \\
\hline
\textbf{تهیه‌کنندگان} & تیم توسعه کلاس یار \\
\hline
\textbf{وضعیت} & پیش‌نویس \\
\hline
\end{tabular}
\end{center}

\newpage
\tableofcontents
\newpage

% ============================================================
\section{مقدمه}

\subsection{هدف سند}
این سند معماری نرم‌افزار پروژهٔ کلاس یار را تعریف می‌کند.
هدف آن ایجاد یک درک مشترک از ساختار سیستم، تصمیم‌های کلیدی و
محدودیت‌ها است.

\subsection{دامنه}
\begin{itemize}
    \item \textbf{در دامنه:} معماری سطح بالا، لایه‌ها، Endpointها، امنیت.
    \item \textbf{خارج از دامنه:} جزئیات پیاده‌سازی، کد.
\end{itemize}

\subsection{ذی‌نفعان}
\begin{center}
\begin{tabular}{|l|l|}
\hline
\textbf{ذی‌نفع} & \textbf{نقش} \\
\hline
تیم توسعه & پیاده‌سازی \\
\hline
استاد درس & ارزیابی \\
\hline
مشتری فرضی & تأیید \\
\hline
\end{tabular}
\end{center}

% ============================================================
\section{نمای کلی معماری}

\subsection{سبک معماری}
\textbf{Monolith + Layered Architecture}

\begin{itemize}
    \item یک برنامهٔ واحد
    \item سه لایهٔ اصلی: Presentation, Business, Data
    \item مناسب تیم کوچک (۳-۴ نفر)
    \item ساده برای شروع، قابل گسترش
\end{itemize}

\subsection{جریان داده}

\begin{center}
کاربر $\rightarrow$ Frontend (React) $\rightarrow$ API (Express) 
$\rightarrow$ Database (PostgreSQL)
\end{center}

% ============================================================
\section{لایه‌ها}

\subsection{Presentation Layer (Frontend)}

\begin{center}
\begin{tabular}{|l|l|}
\hline
\textbf{مورد} & \textbf{انتخاب} \\
\hline
فناوری & React 18 \\
\hline
Routing & React Router \\
\hline
State & Redux Toolkit \\
\hline
UI & Tailwind CSS \\
\hline
Build & Vite \\
\hline
\end{tabular}
\end{center}

\subsection{Business Layer (Backend)}

\begin{center}
\begin{tabular}{|l|l|}
\hline
\textbf{مورد} & \textbf{انتخاب} \\
\hline
فناوری & Node.js 20 + Express 4 \\
\hline
ORM & Prisma \\
\hline
Validation & Zod \\
\hline
Auth & JWT \\
\hline
\end{tabular}
\end{center}

\subsection{Data Layer (Database)}

\begin{center}
\begin{tabular}{|l|l|}
\hline
\textbf{مورد} & \textbf{انتخاب} \\
\hline
فناوری & PostgreSQL 16 \\
\hline
Migration & Prisma Migrate \\
\hline
Backup & روزانه \\
\hline
Hosting & Railway \\
\hline
\end{tabular}
\end{center}

% ============================================================
\section{تصمیم‌های معماری (ADR)}

\subsection{ADR-01: چرا Monolith؟}

\textbf{تصمیم:} معماری Monolith

\textbf{دلیل:}
\begin{itemize}
    \item تیم کوچک (۳-۴ نفر)
    \item پیچیدگی کم
    \item سرعت توسعه بالاتر
    \item استقرار ساده
\end{itemize}

\subsection{ADR-02: چرا React؟}

\textbf{تصمیم:} React برای Frontend

\textbf{دلیل:}
\begin{itemize}
    \item محبوب و مستند
    \item اکوسیستم بزرگ
    \item SPA برای UX بهتر
\end{itemize}

\subsection{ADR-03: چرا PostgreSQL؟}

\textbf{تصمیم:} PostgreSQL برای Database

\textbf{دلیل:}
\begin{itemize}
    \item Relational و قدرتمند
    \item پشتیبانی از JSON
    \item رایگان و اوپن‌سورس
\end{itemize}

% ============================================================
\section{REST API}

\subsection{Endpointهای عمومی (دانشجو و استاد)}

\begin{center}
\begin{tabular}{|l|l|l|}
\hline
\textbf{Method} & \textbf{Endpoint} & \textbf{توضیح} \\
\hline
GET & /api/notes & لیست جزوه‌ها \\
\hline
GET & /api/notes/:id & جزئیات جزوه \\
\hline
GET & /api/professors & لیست استادان \\
\hline
GET & /api/courses & لیست دروس \\
\hline
GET & /api/search & جست‌وجو و فیلتر \\
\hline
\end{tabular}
\end{center}

\subsection{Endpointهای تعامل (نظرات، امتیازات)}

\begin{center}
\begin{tabular}{|l|l|l|}
\hline
\textbf{Method} & \textbf{Endpoint} & \textbf{توضیح} \\
\hline
POST & /api/notes/:id/rate & ثبت امتیاز برای جزوه \\
\hline
POST & /api/notes/:id/comments & ثبت نظر برای جزوه \\
\hline
POST & /api/professors/:id/rate & ثبت امتیاز برای استاد \\
\hline
POST & /api/professors/:id/comments & ثبت نظر برای استاد \\
\hline
\end{tabular}
\end{center}

\subsection{Endpointهای ادمین}

\begin{center}
\begin{tabular}{|l|l|l|}
\hline
\textbf{Method} & \textbf{Endpoint} & \textbf{توضیح} \\
\hline
POST & /api/admin/professors & افزودن استاد \\
\hline
PUT & /api/admin/professors/:id & ویرایش استاد \\
\hline
DELETE & /api/admin/notes/:id & حذف جزوه نامناسب \\
\hline
GET & /api/admin/reports & لیست گزارش‌ها \\
\hline
\end{tabular}
\end{center}

% ============================================================
\section{امنیت}

\subsection{احراز هویت}
\begin{itemize}
    \item \textbf{JWT} برای API
    \item \textbf{bcrypt} برای رمزنگاری گذرواژه
    \item \textbf{Refresh Token} برای تمدید
\end{itemize}

\subsection{مجوزدهی}
\begin{itemize}
    \item \textbf{Role-Based Access Control (RBAC)}
    \item نقش‌ها: Student, Professor, Admin
\end{itemize}

\subsection{امنیت API}
\begin{itemize}
    \item \textbf{HTTPS} اجباری
    \item \textbf{Rate Limiting} (۱۰۰ درخواست در دقیقه)
    \item \textbf{CORS} محدود
    \item \textbf{Input Validation} با Zod
\end{itemize}

% ============================================================
\section{مقیاس‌پذیری}

\subsection{محدودیت‌های فعلی}
\begin{itemize}
    \item Monolith: مقیاس‌پذیری محدود
    \item یک سرور: تک نقطهٔ شکست
\end{itemize}

\subsection{مسیر رشد}
\begin{enumerate}
    \item فاز ۱ (MVP): Monolith + یک سرور
    \item فاز ۲: Load Balancer + چند سرور
    \item فاز ۳: Database Replication
    \item فاز ۴: جدا کردن سرویس‌های پرمصرف
\end{enumerate}

% ============================================================
\section{ریسک‌ها و بدهی فنی}

\begin{center}
\begin{tabular}{|p{3cm}|c|c|p{4cm}|}
\hline
\textbf{ریسک} & \textbf{احتمال} & \textbf{اثر} & \textbf{راهکار} \\
\hline
رشد سریع کاربران & متوسط & بالا & Load Balancer \\
\hline
نفوذ امنیتی & کم & بالا & OWASP \\
\hline
از دست دادن داده & کم & بالا & Backup روزانه \\
\hline
بدهی فنی & بالا & متوسط & Refactor در Sprint 3 \\
\hline
\end{tabular}
\end{center}

% ============================================================
\section{فناوری‌ها}

\begin{center}
\begin{tabular}{|l|p{7cm}|}
\hline
\textbf{لایه} & \textbf{فناوری} \\
\hline
Frontend & React 18 + Vite + Tailwind \\
\hline
Backend & Node.js 20 + Express 4 \\
\hline
Database & PostgreSQL 16 \\
\hline
ORM & Prisma \\
\hline
Auth & JWT + bcrypt \\
\hline
Testing & Jest + Supertest \\
\hline
CI/CD & GitHub Actions \\
\hline
Deploy & Railway \\
\hline
\end{tabular}
\end{center}

% ============================================================
\clearpage
\phantomsection
\addcontentsline{toc}{section}{پیوست الف — تاریخچهٔ تغییرات}
\section*{پیوست الف — تاریخچهٔ تغییرات}

\vspace{0.5cm}

\begin{center}
\begin{tabular}{|c|c|p{5cm}|p{3cm}|}
\hline
\textbf{نسخه} & \textbf{تاریخ} & \textbf{تغییرات} & \textbf{تهیه‌کننده} \\
\hline
۱.۰ & ۲۹/۰۷/۱۴۰۵ & نسخهٔ اولیه & تیم توسعه کلاس یار \\
\hline
\end{tabular}
\end{center}

\end{document}